\section{THIẾT KẾ KHO DỮ LIỆU}

\subsection{Kiến trúc tổng thể}
\hspace{1cm}
Kiến trúc luồng dữ liệu được tổ chức theo mô hình kinh điển của Kimball: dữ liệu thô được
nạp vào vùng \textit{staging}, sau đó biến đổi và nạp vào kho dữ liệu (DW) theo mô hình
\textit{star schema}, từ đó cung cấp cho hai nhóm tiêu thụ: (1) các \textit{OLAP view} phục
vụ truy vấn báo cáo trực tiếp, và (2) tầng đặc trưng (\textit{feature store}) phục vụ các mô
hình DSS (bảng \texttt{FACT\_ORDER\_FULFILLMENT}). Hình~\ref{fig:kientruc} minh họa toàn bộ luồng này.

\begin{figure}[H]
\centering
\begin{tikzpicture}[
    box/.style={draw, rounded corners, minimum width=2.6cm, minimum height=1cm,
                align=center, font=\small, fill=blue!6},
    dwbox/.style={box, fill=blue!15},
    dssbox/.style={box, fill=green!15},
    lbl/.style={font=\scriptsize, align=center},
    arr/.style={-{Stealth[length=2.2mm]}, thick}
]
\node[box, minimum width=2.8cm] (raw) at (1.3,0) {Dữ liệu nguồn\\(Olist, Kaggle)};
\node[box] (staging) at (5.9,0) {Staging\\(PostgreSQL)};
\node[dwbox] (dw) at (10.1,0) {Data Warehouse\\(Star Schema)};
\node[box, minimum width=3cm] (olap) at (14.2,1.1) {OLAP Views\\(Materialized)};
\node[dssbox, minimum width=3cm] (dss) at (14.2,-1.1) {DSS Service\\(Model, :8001)};
% \node[box, minimum width=3cm] (backend) at (14.2,-3.3) {Backend API\\(BFF, :8000)};
% \node[dwbox, minimum width=4.2cm] (fe) at (7.6,-3.3) {Frontend (React, :5173)\\Manager Dashboard · Ops Console\\Campaign · Risk Predictor};

\draw[arr] (raw) -- node[above, lbl]{Extract} (staging);
\draw[arr] (staging) -- node[above, lbl]{Transform\\+ Load} (dw);
\draw[arr] (dw.east) -- (olap.west);
\draw[arr] (dw.east) -- (dss.west);
% \draw[arr] (olap.east) -- ++(0.5,0) |- (backend.east);
% \draw[arr] (dss) -- (backend);
% \draw[arr] (backend) -- node[above, lbl]{REST / JSON} (fe);
% \node[lbl, text=black!60] at (5.9,-1.1) {Điều phối bởi Airflow:\\staging\_load $\to$ dw\_load $\to$ dss\_batch};
\end{tikzpicture}
\caption{Kiến trúc tổng thể: Raw data $\to$ Staging $\to$ DW $\to$ OLAP/DSS $\to$ Backend API $\to$ Frontend}
\label{fig:kientruc}
\end{figure}

Toàn bộ pipeline ETL được điều phối bởi Apache Airflow thông qua ba DAG:
\texttt{staging\_load}, \texttt{dw\_load} và \texttt{dss\_batch}, chạy tuần tự
theo thứ tự (Mục~3.4).

\subsection{Mô hình đa chiều (Star Schema)}
\hspace{1cm}
Kho dữ liệu được thiết kế theo mô hình star schema (dạng \textit{fact constellation}),
gồm hai bảng \textit{fact} dùng chung sáu bảng \textit{dimension}, phù hợp với đặc thù các câu
hỏi phân tích đã xác định ở Mục 2.3. Hình~\ref{fig:starschema} minh họa lược đồ (\texttt{DIM\_DATE} đóng
vai trò \textit{role-playing} cho 5 mốc ngày của đơn; trên hình chỉ vẽ khóa ngày mua
\texttt{date\_sk} / \texttt{purchase\_date\_sk} mà các phân tích đang dùng).

\begin{figure}[H]
\centering
\definecolor{ssbg}{HTML}{F0F4F8}
\definecolor{ssline}{HTML}{7A8794}
\definecolor{ssred}{HTML}{C44E5B}
\definecolor{ssorange}{HTML}{DB8744}
\definecolor{ssgray}{HTML}{7A8794}
\definecolor{ssyellow}{HTML}{F0C44C}
\definecolor{ssteal}{HTML}{4FAE9F}
\definecolor{ssblue}{HTML}{4A86D4}
\definecolor{sspurple}{HTML}{B03FC8}
% \ssdb[a]{tên node}{(x,y)}{màu}{nhãn}{chú thích}{lblbelow|lblright}: khối trụ CSDL; a = bán kính ngang
\newcommand{\ssdb}[7][0.55]{%
  \begin{scope}[shift={#3}]
    \pgfmathsetmacro{\ra}{#1}\pgfmathsetmacro{\rb}{0.32*#1}\pgfmathsetmacro{\hh}{0.78*#1}
    \fill[#4] (-\ra,\hh) -- (-\ra,-\hh) arc[start angle=180,end angle=360,x radius=\ra,y radius=\rb]
              -- (\ra,\hh) arc[start angle=0,end angle=180,x radius=\ra,y radius=\rb] -- cycle;
    \foreach \y in {0.28,-0.28}
      \draw[white, line width=1.1pt] (-\ra,2*\y*\hh) arc[start angle=180,end angle=360,x radius=\ra,y radius=\rb];
    \filldraw[fill=#4, draw=white, line width=1.1pt] (0,\hh) ellipse[x radius=\ra, y radius=\rb];
    \node[minimum width=2*\ra cm, minimum height=2*\hh cm+2*\rb cm, inner sep=0] (#2) at (0,0) {};
  \end{scope}
  \node[#7=#2, font=\small\ttfamily] {#5\\[-1pt]{\scriptsize\rmfamily #6}};
}
\resizebox{\textwidth}{!}{%
\begin{tikzpicture}[
    fk/.style={line width=1.1pt, ssline},
    key/.style={fill=ssbg, inner sep=2pt, font=\footnotesize\ttfamily, text=black!80, align=center},
    lblbelow/.style={below=3pt of #1, align=center},
    lblright/.style={right=6pt of #1, align=left}
]
% --- fact (bên trái)
\ssdb[0.8]{foi}{(-6,1.5)}{ssred}{fact\_order\_item}{grain: 1 sản phẩm trong 1 đơn}{lblbelow}
\ssdb[0.8]{fof}{(-6,-1.5)}{ssorange}{fact\_order\_fulfillment}{grain: 1 đơn (accumulating snapshot)}{lblbelow}
% --- dimension (bên phải)
\foreach \dimn/\y/\c/\n/\note in {date/3.75/sspurple/dim\_date/role-playing,
                               customer/2.25/ssteal/dim\_customer/SCD Type 2 (địa chỉ),
                               location/0.75/ssblue/dim\_location/{city, state, region},
                               product/-0.75/ssyellow/dim\_product/category,
                               seller/-2.25/ssteal/dim\_seller/seller state,
                               payment/-3.75/ssgray/dim\_payment\_type/phương thức thanh toán}
  {\ssdb{\dimn}{(3.6,\y)}{\c}{\n}{\note}{lblright}}
% --- khoá ngoại: 1 trục dọc nối cả 2 fact; từ trục đi ra từng dimension (mũi tên = tham chiếu)
\draw[fk] (-2.6,3.75) -- (-2.6,-3.75);
\draw[fk] (foi.east) -- (-2.6,1.5);
\draw[fk] (fof.east) -- (-2.6,-1.5);
\fill[ssline] (-2.6,1.5) circle (2.2pt) (-2.6,-1.5) circle (2.2pt);
\foreach \dimn/\y/\k in {date/3.75/{date\_sk\\[-1pt]purchase\_date\_sk},
                      customer/2.25/customer\_sk,
                      location/0.75/location\_sk,
                      product/-0.75/product\_sk,
                      seller/-2.25/seller\_sk,
                      payment/-3.75/payment\_type\_sk} {
  \draw[fk, -{Stealth[length=7pt]}] (-2.6,\y) -- node[key] {\k} (\dimn.west);
  \fill[ssline] (-2.6,\y) circle (2.2pt);
}
% --- nền
\begin{scope}[on background layer]
  \fill[ssbg] ($(current bounding box.south west)+(-0.5,-0.4)$)
              rectangle ($(current bounding box.north east)+(0.5,0.4)$);
\end{scope}
\end{tikzpicture}}
\caption{Star schema của kho dữ liệu: hai bảng fact dùng chung sáu bảng
dimension}
\label{fig:starschema}
\end{figure}
\begin{itemize}
\item Truy cập diagram đầy đủ tại đây: 
\href{https://dbdiagram.io/d/star_schema_edss-6af2bffc6fcc7ae4c7a61917}{ Star Schema Diagram.}
\end{itemize}
\subsubsection{Bảng Fact}
\begin{itemize}
    \item \textbf{FACT\_ORDER\_ITEM} 
    \begin{itemize}
      \item Mỗi dòng ứng với một sản phẩm trong một đơn hàng (tương ứng một dòng \texttt{olist\_order\_items}). 
      \item Khóa ngoại đến \texttt{date\_sk}, \texttt{customer\_sk}, \texttt{product\_sk}, \texttt{seller\_sk},
          \texttt{location\_sk}, \texttt{payment\_type\_sk}; chiều suy
          biến (degenerate) \texttt{order\_id}, \texttt{order\_status}; các độ đo (measures):
          \texttt{quantity}, \texttt{price} (doanh thu), \texttt{freight\_value},
          \texttt{review\_score}. 
      %     Bảng còn lưu sẵn \texttt{total\_value}
      %     (= \texttt{price} + \texttt{freight\_value}), \texttt{delivery\_days} và
      %     \texttt{is\_late\_delivery} ở mức dòng sản phẩm, nhưng các phân tích hiện tại không dùng
      %     ba cột này: 
      \item Số liệu giao hàng được lấy từ \texttt{FACT\_ORDER\_FULFILLMENT} (một dòng
          một đơn) để không bị đếm lặp theo số sản phẩm trong đơn.
    \end{itemize}
    \item \textbf{FACT\_ORDER\_FULFILLMENT} 
      \begin{itemize}
            \item Mỗi dòng ứng với một đơn hàng, thiết kế dạng \textit{accumulating snapshot}: 
            mỗi mốc trong vòng đời đơn (mua, duyệt, bàn giao cho đơn vị vận chuyển, giao tới khách, ngày hẹn giao) là một khóa ngoại tới
          \texttt{DIM\_DATE}; dòng được cập nhật khi đơn đi tới mốc tiếp theo. 
            \item Độ đo cấp đơn: \texttt{n\_items}, \texttt{total\_price}, \texttt{total\_freight},
          \texttt{total\_weight\_g}, \texttt{payment\_installments}, \texttt{distance\_km}
          (khoảng cách người bán -- khách tính từ tọa độ zip), \texttt{estimated\_days},
          \texttt{carrier\_days}, \texttt{actual\_days}, \texttt{delay\_days}, cùng hai nhãn
          \texttt{is\_late} và \texttt{is\_low\_review} (review $\leq$ 2 sao). Đây là nguồn dữ
          liệu của hai mô hình DSS (Mục~\ref{sec:dss}).
      \end{itemize}
\end{itemize}

\subsubsection{Bảng Dimension}
\begin{table}[H]
\centering
\caption{Danh sách bảng chiều (dimension) và chiến lược SCD}
\begin{tabular}{@{}p{4cm}p{2cm}p{10cm}@{}}
\toprule
\textbf{Bảng} & \textbf{SCD} & \textbf{Ghi chú} \\
\midrule
\texttt{DIM\_CUSTOMER} & Type 2 & Khóa nghiệp vụ \texttt{customer\_unique\_id}. Lưu lịch sử
thay đổi địa chỉ (zip, thành phố, bang) qua các cột \texttt{effective\_from},
\texttt{effective\_to} (kiểu \texttt{TIMESTAMP}) và \texttt{is\_current}; 252 khách có hơn
một phiên bản. \\
\texttt{DIM\_PRODUCT} & Type 1 & Mã sản phẩm và ngành hàng (đã dịch sang tiếng Anh); ghi đè
khi thay đổi. \\
\texttt{DIM\_DATE} & Static & Sinh sẵn cho các năm có dữ liệu, gồm ngày, thứ trong tuần, tháng,
năm. \\
\texttt{REF\_ZIP\_GEO} & Tham chiếu & Tọa độ trung bình của 19.010 zip prefix (từ
\texttt{olist\_geolocation}, đã bỏ điểm nằm ngoài Brazil), dùng tính \texttt{distance\_km}. \\
\texttt{DIM\_LOCATION} & Type 1 & Thành phố, bang, vùng địa lý (nơi giao hàng). \\
\texttt{DIM\_SELLER} & Type 1 & Người bán: mã, bang. \\
\texttt{DIM\_PAYMENT\_\allowbreak TYPE} & Type 1 & Phương thức thanh toán chính của đơn (credit card,
boleto, voucher, debit card). Dùng thay cho chiều ``kênh bán'' vì Olist không có thông tin
kênh. \\
\bottomrule
\end{tabular}
\end{table}

Các bảng chiều chỉ giữ những thuộc tính đang được OLAP view, API hoặc mô hình DSS truy vấn,
cùng các khóa tự nhiên cần cho ETL; dựa trên mức độ cần thiết đối với dự án hiện tại, một số thuộc tính được lược bỏ 
nhưng vẫn có giá trị và khả năng phát triển vào các giai đoạn sau của dự án (Mục~\ref{sec:dw-gaps}).

% Việc chọn \textbf{SCD Type 2} cho \texttt{DIM\_CUSTOMER} xuất phát từ đặc điểm dữ liệu Olist:
% địa chỉ giao hàng được ghi theo từng đơn, và một số khách chuyển nơi nhận hàng giữa các lần
% mua. Với SCD Type 1, các đơn cũ sẽ bị gán sang địa chỉ mới, làm sai thống kê theo khu vực
% (ví dụ tỷ lệ giao trễ theo bang). Với SCD Type 2, bảng fact được nối theo thời điểm mua
% (\textit{point-in-time join}: \texttt{effective\_from} $\leq$ thời điểm mua $<$
% \texttt{effective\_to}) nên mỗi đơn luôn trỏ tới đúng phiên bản địa chỉ tại thời điểm đặt hàng.
% Mốc hiệu lực được lưu ở mức \textbf{thời điểm} chứ không phải ngày: trong dữ liệu có 14 khách
% đặt hai đơn tới hai địa chỉ khác nhau trong cùng một ngày; khi mốc chỉ tính theo ngày, 12 đơn bị
% nối sai phiên bản (5 đơn sai cả bang). Lỗi này được phát hiện khi đối soát fact với dimension và
% đã được sửa, kèm một kiểm tra tự động (Mục~3.5).

\subsubsection{Bảng vận hành cho chiến dịch chăm sóc}
Hai bảng \texttt{dw.care\_campaign} và \texttt{dw.campaign\_target} không thuộc star schema.
Đây là bảng lưu trữ phục vụ DSS, do người dùng thay đổi trực tiếp trên giao diện
(\textit{write-back} qua backend API); ETL không đọc và không ghi vào hai bảng này:
\begin{itemize}
    \item \texttt{care\_campaign}: mã, tên, mô tả chiến dịch do Ops tạo trên giao diện. Thao tác
          xóa là \textit{xóa mềm} (\texttt{deleted\_at}) để giữ lịch sử chiến dịch.
    \item \texttt{campaign\_target}: khách hàng đã gắn vào từng chiến dịch, kèm đơn và điểm rủi
          ro dẫn tới việc chọn khách; mỗi khách xuất hiện tối đa một lần trong một chiến dịch.
\end{itemize}

\subsection{OLAP Views}
\hspace{1cm}
Để tránh việc tính toán lại các phép tổng hợp phức tạp mỗi lần truy vấn dashboard, ba
\textit{materialized view} được xây dựng trực tiếp trên kho dữ liệu:

\begin{table}[H]
\centering
\caption{OLAP views phục vụ báo cáo và DSS}
\begin{tabular}{@{}p{7.0cm}>{\raggedright\arraybackslash}p{9.4cm}@{}}
\toprule
\textbf{View} & \textbf{Mục đích} \\
\midrule
\texttt{mv\_revenue\_by\_\allowbreak month\_category} & Rollup doanh thu theo tháng và ngành hàng --
trả lời câu hỏi nghiệp vụ (1). \\
\texttt{mv\_top\_products} & Xếp hạng sản phẩm bán chạy/chậm theo doanh thu -- trả lời câu
hỏi nghiệp vụ (2). \\
\texttt{mv\_delivery\_performance} & Đơn đã giao theo tháng mua $\times$ bang người bán
$\times$ vùng/bang khách: số đơn, số đơn trễ, tổng số ngày giao/ngày hẹn, tổng điểm review,
số review thấp (lưu tổng thay vì trung bình để rollup đúng) -- trả lời câu hỏi nghiệp vụ (3). \\
\bottomrule
\end{tabular}
\end{table}

\subsection{Quy trình ETL}
\hspace{1cm}
Quy trình ETL được cài đặt bằng Python và điều phối bởi Airflow, gồm ba giai đoạn tương ứng
ba module trong \texttt{etl/}:

\begin{enumerate}
    \item \textbf{Extract} (\texttt{etl/extract/extract\_olist.py}): đọc 9 file CSV Olist ($\sim$1 triệu dòng)
          , chỉ chọn các cột cần dùng, \texttt{TRUNCATE} rồi
          \texttt{COPY} vào schema \texttt{staging} của PostgreSQL.
      %   (full reload, chạy lại không nhân đôi dữ liệu).
    \item \textbf{Transform} (\texttt{etl/transform/transform\_dw.py}): xây dựng các bảng
          dimension bằng SQL set-based (\texttt{INSERT ... ON CONFLICT}): sinh
          \texttt{DIM\_DATE}, gom tọa độ theo zip vào \texttt{REF\_ZIP\_GEO}, dịch tên ngành hàng
          sang tiếng Anh, ánh xạ bang sang vùng địa lý, gom khách hàng theo
          \texttt{customer\_unique\_id} và áp dụng SCD Type 2 theo địa chỉ cho
          \texttt{DIM\_CUSTOMER}.
    \item \textbf{Load} (\texttt{etl/load/load\_dw.py}): nạp hai bảng fact, tra khóa,
      %     (với \texttt{DIM\_CUSTOMER} lấy phiên bản có hiệu lực tại ngày mua)
          tính các độ đo dẫn xuất (\texttt{total\_value}, \texttt{delivery\_days}, \texttt{is\_late\_delivery},
          điểm review trung bình của đơn; với \texttt{FACT\_ORDER\_FULFILLMENT} thêm khoảng cách
          haversine, các khoảng thời gian giữa các mốc và nhãn trễ/review thấp). Cả hai fact đều
          dùng \texttt{ON CONFLICT DO UPDATE} nên khi chạy lại, trạng thái đơn, ngày giao, nhãn
          trễ và điểm review được cập nhật theo vòng đời đơn thay vì giữ giá trị cũ; sau đó \texttt{REFRESH MATERIALIZED VIEW} cho các
          OLAP view liên quan.
\end{enumerate}

Ba Airflow DAG (thư mục \texttt{etl/dags/}) nối tiếp nhau bằng
\texttt{TriggerDagRunOperator}:
\begin{itemize}
    \item \texttt{staging\_load}: chạy Extract, nạp dữ liệu thô vào staging. Chạy thủ công
          (\texttt{schedule=None}) vì Olist là dataset lịch sử tĩnh.
    \item \texttt{dw\_load}: chạy Transform + Load, cập nhật kho dữ liệu và OLAP view,
          kích hoạt sau khi \texttt{staging\_load} hoàn tất.
    \item \texttt{dss\_batch}: huấn luyện lại / chạy suy luận theo lô (batch inference)
          cho ba mô hình DSS, kích hoạt sau khi \texttt{dw\_load} hoàn tất (và chạy định kỳ
          \texttt{@weekly}).
\end{itemize}

\subsection{Kiểm tra chất lượng dữ liệu}
\hspace{1cm}
Trước khi dữ liệu từ staging được nạp vào kho, các kiểm tra chất lượng sau được thực hiện
(cài đặt tại \texttt{etl/tests/test\_data\_quality.py}):
\begin{itemize}
    \item Kiểm tra giá trị NULL trên các khóa nghiệp vụ bắt buộc (\texttt{order\_id},
          \texttt{customer\_id}, \texttt{product\_id}, ...) và không trùng
          (\texttt{order\_id}, \texttt{order\_item\_id}).
    \item Kiểm tra ràng buộc toàn vẹn tham chiếu (referential integrity): mọi dòng
          \texttt{order\_items} trỏ tới đơn hàng hợp lệ, mọi đơn trỏ tới khách hàng hợp lệ;
          toàn bộ \texttt{order\_items} đều được nạp vào fact và có đủ
          \texttt{location\_sk}, \texttt{payment\_type\_sk}.
    \item Kiểm tra miền giá trị hợp lệ: \texttt{price} $\geq$ 0, \texttt{review\_score}
          trong khoảng 1--5.
    \item Kiểm tra SCD2: mỗi khách hàng có đúng một phiên bản hiện hành, các phiên bản không
          chồng lấn, \texttt{DIM\_CUSTOMER} dùng \texttt{customer\_unique\_id} làm khóa
          nghiệp vụ, và mỗi dòng fact trỏ tới phiên bản có đúng địa chỉ giao của đơn.
    \item Kiểm tra \texttt{FACT\_ORDER\_FULFILLMENT}: đủ một dòng cho mỗi đơn có sản phẩm, nhãn
          \texttt{is\_late}/\texttt{is\_low\_review} khớp với \texttt{delay\_days}/\texttt{review\_score},
          mọi đơn đều có \texttt{distance\_km} (16 kiểm tra, tất cả đạt).
\end{itemize}
Bộ kiểm tra viết bằng \texttt{pytest}, chạy ở bước cuối của \texttt{scripts/seed\_db.sh}
(\texttt{pytest etl/tests}); nếu có kiểm tra thất bại, script dừng lại để nhóm xử lý trước
khi đưa dữ liệu vào phân tích.
